% Part: computability
% Chapter: computability-theory
% Section: rice-theorem

\documentclass[../../../include/open-logic-section]{subfiles}

\begin{document}

\olfileid{cmp}{thy}{rce}
\olsection{রাইসের উপপাদ্য}

একটু ভাবলে দেখা যাবে, \olref[tot]{prop:total}-এর প্রমাণে $\fn{Tot}$-এর
বিশেষ বৈশিষ্ট্য কোনও ভূমিকা নেয় না। ঠিক যখন $x$, $K$-তে থাকে তখনই
$h(x,y)$ যেন ধ্রুব অপেক্ষক $j(y) = 0$-এর মতো আচরণ করে, সেভাবে আমরা তাকে
নির্মাণ করেছি; কিন্তু সেই পরিস্থিতিতে তাকে অন্য যেকোনো আংশিক গণনসাধ্য
অপেক্ষকের মতো আচরণ করানোও যেত। এই পর্যবেক্ষণটি একটি আরও সাধারণ উপপাদ্য
বলার সুযোগ দেয়: মোটামুটি বললে, গণনসাধ্য অপেক্ষকের কোনো অতুচ্ছ বৈশিষ্ট্যই
নির্ণেয় নয়।

মনে রেখো, $\cfind{0}$, $\cfind{1}$, $\cfind{2}$,~\dots হলো আংশিক
গণনসাধ্য অপেক্ষকগুলির আমাদের নির্দিষ্ট তালিকায়ন।

\begin{thm}[রাইসের উপপাদ্য]
  $C$-কে আংশিক গণনসাধ্য অপেক্ষকগুলির যেকোনো সেট ধরি এবং
  $A = \Setabs{n}{\cfind{n} \in C}$ ধরি। $A$ গণনসাধ্য হলে হয় $C$
  শূন্য, নয়তো $C$ সব আংশিক গণনসাধ্য অপেক্ষকের সেট।
\end{thm}

একটি \emph{সূচক-সেট} হলো এমন একটি সেট~$A$, যার বৈশিষ্ট্য এই: $n$ ও $m$
যদি একই অপেক্ষক “গণনা” করে এমন সূচক হয়, তবে হয় $n$ ও $m$ উভয়েই
$A$-তে থাকে, নয়তো কেউই থাকে না। উপপাদ্যের সেট~$A$-এর এই বৈশিষ্ট্য আছে
তা বোঝা কঠিন নয়। বিপরীতভাবে, $A$ একটি সূচক-সেট এবং $C$ এই সূচকগুলি
দিয়ে গণিত অপেক্ষকগুলির সেট হলে $A = \Setabs{n}{\cfind{n} \in C}$।

\begin{explain}
এই পরিভাষায় রাইসের উপপাদ্যটি এই বক্তব্যের সমতুল্য যে কোনো অতুচ্ছ
সূচক-সেট নির্ণেয় নয়। উপপাদ্যটি কী বলে তা বুঝতে \emph{প্রোগ্রাম}
(ধরা যাক, তোমার পছন্দের প্রোগ্রামিং ভাষায়) এবং প্রোগ্রাম যে অপেক্ষক
গণনা করে—এ দুটির পার্থক্যে জোর দেওয়া সহায়ক। প্রোগ্রাম (সূচক) হলো
বাক্যগঠনগত বস্তু, এবং তাদের সম্পর্কে অবশ্যই গণনসাধ্য প্রশ্ন আছে: এই
প্রোগ্রামে কি 150টির বেশি প্রতীক আছে? এতে কি 22টির বেশি পঙ্‌ক্তি আছে?
এতে কি কোনো “while” বিবৃতি আছে? কোনো “print” বিবৃতির আর্গুমেন্ট হিসেবে
“hello world” প্রতীকক্রমটি কি কখনও আসে? রাইসের উপপাদ্য বলে, প্রোগ্রামের
\emph{আচরণ} সম্পর্কে কোনো অতুচ্ছ প্রশ্ন গণনসাধ্য নয়। এর মধ্যে এমন
প্রশ্নও পড়ে: প্রোগ্রামটি কি নিবেশ~$0$-এ থামে? সেটি কি কখনও থামে? সেটি
কি কখনও কোনো জোড় সংখ্যা নির্গত করে?
\end{explain}

\begin{proof}[রাইসের উপপাদ্যের প্রমাণ]
ধরা যাক, $C$ শূন্যও নয়, সব আংশিক গণনসাধ্য অপেক্ষকের সেটও নয়, এবং
$A$ হলো~$C$-এর অপেক্ষকগুলির সূচক-সেট। আমরা দেখাব, $A$ গণনসাধ্য হলে
থামার সমস্যার সমাধান করা যেত; তাই~$A$ গণনসাধ্য নয়।

সাধারণত্ব না হারিয়ে ধরে নিতে পারি, সর্বত্র অসংজ্ঞায়িত অপেক্ষক~$f$,
$C$-তে নেই (অন্যথায় নিচের যুক্তিতে $C$ ও তার পূরক অদলবদল করো)।
$g$-কে~$C$-এর যেকোনো অপেক্ষক ধরি। ধারণাটি হলো, $A$-কে নির্ণয় করা গেলে
$f$ গণনাকারী সূচক ও $g$ গণনাকারী সূচকের পার্থক্য বলা যেত; এরপর সেই
সামর্থ্য দিয়ে থামার সমস্যার সমাধান করা যেত।

কীভাবে, তা দেখা যাক। সার্বজনীন আংশিক গণনসাধ্য অপেক্ষক ব্যবহার করে আমরা
একটি অপেক্ষক সংজ্ঞায়িত করতে পারি:
\[
h(x,y) \simeq
\begin{cases}
\text{অসংজ্ঞায়িত} & \text{যদি $\cfind{x}(x) \fundefined$} \\
g(y) & \text{অন্যথায়।}
\end{cases}
\]
$h$ গণনা করতে প্রথমে $\cfind{x}(x)$ গণনা করার চেষ্টা করি; সেই গণনা
থামলে $g(y)$ গণনা করতে এগোই; এবং \emph{সেই} গণনা থামলে নির্গমটি ফেরত
দিই। আরও আনুষ্ঠানিকভাবে লেখা যায়,
\[
h(x,y) \simeq \Proj{2}{0}(g(y),\fn{Un}(x,x)).
\]
এখানে $\Proj{2}{0}(z_0, z_1) = z_0$ হলো $0$-তম আর্গুমেন্ট ফেরত দেয়
এমন গণনসাধ্য $2$-স্থানী অভিক্ষেপ অপেক্ষক।

তাহলে $h$ আংশিক গণনসাধ্য অপেক্ষকগুলির একটি সংযোজন, এবং ডান পাশটি ঠিক
তখনই সংজ্ঞায়িত ও~$g(y)$-এর সমান, যখন $\fn{Un}(x,x)$ ও $g(y)$ উভয়ই
সংজ্ঞায়িত।

লক্ষ করো, কোনো স্থির~$x$-এর জন্য $\cfind{x}(x)$ অসংজ্ঞায়িত হলে প্রত্যেক
~$y$-এর জন্য $h(x,y)$ অসংজ্ঞায়িত; আর $\cfind{x}(x)$ সংজ্ঞায়িত হলে
$h(x,y) \simeq g(y)$। তাই~$x$-এর যেকোনো স্থির মানের জন্য $h(x,y)$ হয়
ঠিক~$f$-এর মতো, নয়তো ঠিক~$g$-এর মতো আচরণ করে; $\cfind{x}(x)$
সংজ্ঞায়িত কি না নির্ণয় করা মানে এই দুই ক্ষেত্রের কোনটি ঘটে তা নির্ণয়
করা। কিন্তু এর অর্থ, $h_x(y) \simeq h(x,y)$, $C$-তে আছে কি না নির্ণয়
করা; $A$ গণনসাধ্য হলে আমরা ঠিক সেটিই করতে পারতাম।

আরও আনুষ্ঠানিকভাবে, $h$ আংশিক গণনসাধ্য বলে কোনো সূচক~$e$-এর জন্য সেটি
$\cfind{e}$-এর সমান। $s$-$m$-$n$ উপপাদ্য অনুসারে এমন একটি আদিম
পুনরাবৃত্ত অপেক্ষক~$s$ আছে, যাতে প্রত্যেক~$x$-এর জন্য
$\cfind{s(e,x)}(y) \simeq h_x(y)$। এখন প্রত্যেক~$x$-এর জন্য
$\cfind{x}(x) \fdefined$ হলে $\cfind{s(e,x)}$, $g$-এর একই অপেক্ষক, তাই
$s(e,x)$, $A$-তে আছে। অন্যদিকে, $\cfind{x}(x) \uparrow$ হলে
$\cfind{s(e,x)}$, $f$-এর একই অপেক্ষক, তাই $s(e,x)$, $A$-তে নেই। অন্যভাবে
বললে, প্রত্যেক~$x$-এর জন্য $x \in K$ যদি এবং কেবল যদি
$s(e,x) \in A$। $A$ গণনসাধ্য হলে~$K$-ও তাই হতো, যা বিরোধ। সুতরাং $A$
গণনসাধ্য নয়।
% BN-SRC-156 TARGET-CORRECTION-BEGIN
\par\smallskip\noindent\textnormal{\small\textbf{উৎস-সংশোধন (BN-SRC-156):} হিমায়িত ইংরেজি প্রমাণে বিশেষায়িত আংশিক অপেক্ষকের সমীকরণে সাধারণ “$=$” আছে। উভয় পাশ অসংজ্ঞায়িত হতে পারে বলে বাংলা পাঠে যুগপৎ সংজ্ঞায়িততার সমতা “$\simeq$” রাখা হয়েছে।} % BN-SRC-156 TARGET-CORRECTION-END
\end{proof}

রাইসের উপপাদ্য অত্যন্ত শক্তিশালী। নিচের তাৎক্ষণিক অনুসিদ্ধান্তে কয়েকটি
নমুনা প্রয়োগ দেখা যায়।
\begin{cor}
নিচের সেটগুলি অনির্ণেয়।
\begin{enumerate}
\item $\Setabs{x}{\text{$17$, $\cfind{x}$-এর মানপরিসরে আছে}}$
\item $\Setabs{x}{\text{$\cfind{x}$ ধ্রুব}}$
\item $\Setabs{x}{\text{$\cfind{x}$ সর্বত্র-সংজ্ঞায়িত}}$
\item $\Setabs{x}{\text{যখনই $y < y'$, এবং $\cfind{x}(y) \fdefined$ ও
    $\cfind{x}(y') \fdefined$, তখন $\cfind{x}(y) < \cfind{x}(y')$}}$
\end{enumerate}
% BN-SRC-157 TARGET-CORRECTION-BEGIN
\par\smallskip\noindent\textnormal{\small\textbf{উৎস-সংশোধন (BN-SRC-157):} হিমায়িত ইংরেজি চতুর্থ দফায় “whenever” ধারার পরে $\cfind{x}(y)$-এর সংজ্ঞায়িততা আলাদা শর্ত হিসেবে বসিয়ে আবার শুধু $\cfind{x}(y')$-এর আগে “if” দেয়, ফলে বাক্যটির যৌক্তিক গঠন অসম্পূর্ণ। বাংলা পাঠে উদ্দেশ্য অনুযায়ী বলা হয়েছে: $y<y'$ এবং উভয় মান সংজ্ঞায়িত হলে প্রথম মানটি দ্বিতীয়টির চেয়ে ছোট।} % BN-SRC-157 TARGET-CORRECTION-END
\end{cor}

\begin{proof}
  এগুলির সবকটিই অতুচ্ছ সূচক-সেট।
\end{proof}

\end{document}
